\documentclass[10pt,a4paper]{amsart}
\usepackage{amsmath,amssymb,amsthm,graphicx}
\usepackage{verbatim}
\newenvironment{reference}{\comment}{\endcomment}
%\newenvironment{reference}{}{}
\newenvironment{slogan}{\comment}{\endcomment}
\newenvironment{history}{\comment}{\endcomment}

% For commutative diagrams we use Xy-pic
\usepackage[all]{xy}

% We use 2cell for 2-commutative diagrams.
\xyoption{2cell}
% Theorem environments.
%
\theoremstyle{plain}
\newtheorem{theorem}[subsection]{Theorem}
\newtheorem{proposition}[subsection]{Proposition}
\newtheorem{lemma}[subsection]{Lemma}

\theoremstyle{definition}
\newtheorem{definition}[subsection]{Definition}
\newtheorem{example}[subsection]{Example}
\newtheorem{exercise}[subsection]{Exercise}
\newtheorem{situation}[subsection]{Situation}

\theoremstyle{remark}
\newtheorem{remark}[subsection]{Remark}
\newtheorem{remarks}[subsection]{Remarks}

\numberwithin{equation}{subsection}

% Macros
%
\def\lim{\mathop{\mathrm{lim}}\nolimits}
\def\colim{\mathop{\mathrm{colim}}\nolimits}
\def\Spec{\mathop{\mathrm{Spec}}}
\def\Hom{\mathop{\mathrm{Hom}}\nolimits}
\def\Ext{\mathop{\mathrm{Ext}}\nolimits}
\def\SheafHom{\mathop{\mathcal{H}\!\mathit{om}}\nolimits}
\def\SheafExt{\mathop{\mathcal{E}\!\mathit{xt}}\nolimits}
\def\Sch{\mathit{Sch}}
\def\Mor{\mathop{\mathrm{Mor}}\nolimits}
\def\Ob{\mathop{\mathrm{Ob}}\nolimits}
\def\Sh{\mathop{\mathit{Sh}}\nolimits}
\def\NL{\mathop{N\!L}\nolimits}
\def\CH{\mathop{\mathrm{CH}}\nolimits}
\def\proetale{{pro\text{-}\acute{e}tale}}
\def\etale{{\acute{e}tale}}
\def\QCoh{\mathit{QCoh}}
\def\Ker{\mathop{\mathrm{Ker}}}
\def\Im{\mathop{\mathrm{Im}}}
\def\Coker{\mathop{\mathrm{Coker}}}
\def\Coim{\mathop{\mathrm{Coim}}}

% Boxtimes
%
\DeclareMathSymbol{\boxtimes}{\mathbin}{AMSa}{"02}

%
% Macros for moduli stacks/spaces
%
\def\QCohstack{\mathcal{QC}\!\mathit{oh}}
\def\Cohstack{\mathcal{C}\!\mathit{oh}}
\def\Spacesstack{\mathcal{S}\!\mathit{paces}}
\def\Quotfunctor{\mathrm{Quot}}
\def\Hilbfunctor{\mathrm{Hilb}}
\def\Curvesstack{\mathcal{C}\!\mathit{urves}}
\def\Polarizedstack{\mathcal{P}\!\mathit{olarized}}
\def\Complexesstack{\mathcal{C}\!\mathit{omplexes}}
% \Pic is the operator that assigns to X its picard group, usage \Pic(X)
% \Picardstack_{X/B} denotes the Picard stack of X over B
% \Picardfunctor_{X/B} denotes the Picard functor of X over B
\def\Pic{\mathop{\mathrm{Pic}}\nolimits}
\def\Picardstack{\mathcal{P}\!\mathit{ic}}
\def\Picardfunctor{\mathrm{Pic}}
\def\Deformationcategory{\mathcal{D}\!\mathit{ef}}
\usepackage{cleveref}
\crefname{originallemma}{original lemma}{original lemmas}
\Crefname{originallemma}{Original lemma}{Original lemmas}
\makeatletter
\newlabel{lemma-ext-modules-hom}{{0DVE}{source}{}{sourcepointer-1}{}}
\newlabel{lemma-ext-modules-hom@cref}{{[originallemma][0][]0DVE}{[1][0][]source}{}{}{}}
\makeatother
\crefname{originaldefinition}{original definition}{original definitions}
\Crefname{originaldefinition}{Original definition}{Original definitions}
\makeatletter
\newlabel{definition-G-module-continuous}{{04JP}{source}{}{sourcepointer-2}{}}
\newlabel{definition-G-module-continuous@cref}{{[originaldefinition][0][]04JP}{[1][0][]source}{}{}{}}
\makeatother
\title{First cohomology of a finite dimensional continuous module over a topologically finitely generated group is finite dimensional}
\author{The Stacks Project Authors}
\date{stacks-project:etale-cohomology.tex:local-snapshot-20261011; GNU FDL 1.2 or later}
\begin{document}
\maketitle
\noindent\textit{Editorial note.} All delivered mathematical text is the author's own, copied from the pinned Stacks Project chapter file etale-cohomology.tex; no statement or proof step was rewritten, reordered or removed.\par\medskip
\noindent\textit{Editorial note.} The proof is the author's complete argument; the parenthetical 'some details omitted' concerns the routine passage to a subquotient of a finite power, and the decisive step that the vectors determine the extension is written out in full.\par\medskip
\noindent\textit{Editorial note.} Every labelled reference occurring inside the counted spans is carried in assembly.external\_references as a non-load-bearing navigation pointer (same-chapter ones at their definition site, cross-chapter ones at the invocation inside the counted span), so no load-bearing mathematics is replaced by a stub. No bibliography citation occurs inside the counted spans.\par\medskip
\noindent\textit{Editorial note.} The checked chapter list for same-name collisions is the set of stacks chapters staged in this repository: chow, cohomology, divisors, etale-cohomology, flat, more-algebra, more-morphisms, morphisms, spaces-more-morphisms, varieties together with this card's chapter etale-cohomology.\par\medskip
\noindent\textit{Editorial note.} The preamble is the chapter's own native preamble excerpt already used for this source: the Xy-pic setup (native preamble lines 14-18) and the native theorem declarations with the mathematical macros (native preamble lines 154-217).\par\medskip
\noindent\textit{Allowed standard prerequisites (editorial).}\begin{itemize}
\item First cohomology as extensions. Same chapter, Lemma lemma-ext-modules-hom: the first cohomology of a discrete module is the space of extensions of the trivial module by it. Turns the group to be bounded into a space of extensions.
\item Continuous modules over a topological group. Same chapter, Definition definition-G-module-continuous. Supplies the continuity condition used to find an open subgroup acting trivially on a chosen lift.
\end{itemize}\medskip
\ifdefined\hypertarget\hypertarget{sourcepointer-1}{}\fi
\noindent\textit{Source reference.} Original lemma 0DVE: sources/stacks-etale-cohomology/pinned-original-source.tex, source lines 8906--8906. Navigation pointer invoked at proof line 8951; the counted span contains the pointer only and the cited result is not reproduced inside it, so the pointer is registered as a non-load-bearing navigation stub in the established form for Stacks material. No bibliography citation is involved.\par\medskip
\ifdefined\hypertarget\hypertarget{sourcepointer-2}{}\fi
\noindent\textit{Source reference.} Original definition 04JP: sources/stacks-etale-cohomology/pinned-original-source.tex, source lines 8786--8786. Navigation pointer invoked at proof line 8968; the counted span contains the pointer only and the cited result is not reproduced inside it, so the pointer is registered as a non-load-bearing navigation stub in the established form for Stacks material. No bibliography citation is involved.\par\medskip
\begin{lemma}
\label{lemma-finite-dim-group-cohomology}
Let $G$ be a topological group. Let $k$ be a field.
Let $V$ be a $k\text{-}G$-module.
If $G$ is topologically finitely generated and
$\dim_k(V) < \infty$, then $\dim_k H^1(G, V) < \infty$.
\end{lemma}

\begin{proof}
Let $g_1, \ldots, g_r \in G$ be elements which topologically generate $G$,
i.e., this means that the subgroup generated by $g_1, \ldots, g_r$ is dense.
By Lemma \ref{lemma-ext-modules-hom}
we see that $H^1(G, V)$ is the $k$-vector space of extensions
$$
0 \to V \to E \to k \to 0
$$
of $k\text{-}G$-modules. Choose $e \in E$ mapping to $1 \in k$.
Write
$$
g_i \cdot e = v_i + e
$$
for some $v_i \in V$. This is possible because $g_i \cdot 1 = 1$.
We claim that the list of elements $v_1, \ldots, v_r \in V$
determine the isomorphism class of the extension $E$.
Once we prove this the lemma follows as this means that our
Ext vector space is isomorphic to a subquotient of the $k$-vector
space $V^{\oplus r}$; some details omitted.
Since $E$ is an object of the category defined in
Definition \ref{definition-G-module-continuous}
we know there is an open subgroup $U$ such that
$u \cdot e = e$ for all $u \in U$.
Now pick any $g \in G$. Then $gU$ contains a word $w$ in
the elements $g_1, \ldots, g_r$.
Say $gu = w$. Since the element $w \cdot e$ is determined by
$v_1, \ldots, v_r$, we see that $g \cdot e = (gu) \cdot e = w \cdot e$
is too.
\end{proof}

\end{document}
